\documentclass[a4paper]{article}

\usepackage{anyfontsize}
\usepackage{amsmath,amssymb,mathrsfs}
\usepackage{autobreak}
\allowdisplaybreaks
\usepackage{bm}
\usepackage{indentfirst}
\setlength{\parindent}{0em}
\usepackage{parskip}
\usepackage{geometry}
\geometry{top=3cm,bottom=3cm,left=2.75cm,right=2.75cm}

\begin{document}

\title{\textbf{Exercise 4}}
\author{}
\date{}

\maketitle

\subsection*{Problem 1: Galaxy groups}

Assume a flat concordance cosmology (i.e. $\Omega_\mathrm{m}=0.27$, $\Omega_\Lambda=0.73$) in the following.

\begin{itemize}
    \item[(a)] Peculiar motions of galaxies usually have amplitudes of a few hundred $\mathrm{km/s}$. What fractional error (i.e. $\delta H_0/H_0$) is made in Hubble's parameter if the distance is measured to a galaxy at cosmological redshift $z$ that has a peculiar velocity of $300\,\mathrm{km/s}$ away from us? Consider the redshifts $z=0.001$, $z=0.01$, and $z=0.1$?
    
    \item[(b)] Groups of galaxies usually have a velocity dispersion of about $500\,\mathrm{km/s}$. If you try to identify such groups in a redshift survey, in which you know the angular positions and redshifts of large numbers of galaxies, would you expect it to be harder or easier to distinguish members of such groups at high redshift compared with groups in the local universe, assuming they have the same velocity dispersions, in terms of confusing unrelated foreground or background galaxies as members?
    
    \item[(c)] The groups in (b) represent haloes of a particular mass $M$, say $10^{14} M_\odot$. They will form over an extended period of time since their number density in the universe must increase quite rapidly with epoch, as desribed by the Press-Schechter analysis. For groups of a given mass $M$, how would you expect the velocity dispersion of those groups that have "just recently virialized" to depend on redshift? Does this consideration alter your answer to (b)?
\end{itemize}

\subsection*{Problem 2: Alcock-Paczy\'{n}ski test}

The so-called “Alcock-Paczy\'{n}ski test” seeks to determine cosmological parameters by means of comparing the radial and angular extents (i.e. $\Delta z$ and $\Delta\theta$) of something that can be assumed to be spherically symmetric. One possibility could be the correlation function of galaxies $\xi(r)$ which, from isotropy, should be a function of the separation $r$ only.

\begin{itemize}
    \item[(a)] Why does this test work in principle?
    
    \item[(b)] How strong is this effect at redshift $z=0.5$? Compare for instance the ratio $\Delta z/\Delta\theta$ for a given object in the concordance cosmology to that in the Einstein-de Sitter cosmology (i.e. $\Omega_\mathrm{m}=1$).
    
    \item[(c)] What complication has made it extremely hard to apply this test in practice?
\end{itemize}

\end{document}